% lpv-lfr-baseline.tex
% --------------------
% Comprehensive documentation of the LPV-LFR baseline implementation
% for the dual-gantry first-principles model.
%
% Compile: pdflatex lpv-lfr-baseline.tex  (twice, for cross-references)
%
\documentclass[11pt,a4paper]{article}

\usepackage[utf8]{inputenc}
\usepackage[T1]{fontenc}
\usepackage{amsmath,amssymb,bm}
\usepackage{booktabs}
\usepackage{geometry}
\usepackage{hyperref}
\usepackage{cleveref}
\usepackage{xcolor}
\usepackage{listings}
\usepackage{enumitem}
\usepackage{caption}

\geometry{margin=2.5cm}
\hypersetup{colorlinks=true, linkcolor=blue!70!black, citecolor=green!50!black, urlcolor=blue!60!black}

% Shorthand commands
\newcommand{\R}{\mathbb{R}}
\newcommand{\inv}{^{-1}}
\newcommand{\T}{^{\!\top}}
\newcommand{\dt}{\Delta t}
\newcommand{\ts}{t_s}
\newcommand{\fs}{f_s}
\newcommand{\xdot}{\dot{x}}
\newcommand{\qdot}{\dot{q}}
\newcommand{\qddot}{\ddot{q}}
\newcommand{\mat}[1]{\mathbf{#1}}
\newcommand{\decref}[1]{\texttt{D-#1}}
\newcommand{\pyfile}[1]{\texttt{#1}}

\lstset{
  basicstyle=\ttfamily\small,
  keywordstyle=\color{blue!70!black},
  commentstyle=\color{green!50!black},
  breaklines=true,
  frame=single,
  columns=fullflexible,
}

\title{LPV-LFR Baseline Model for the ASMPT Dual-Gantry:\\
Implementation, Design Choices, and Verification}
\author{Dirk van den Berg}
\date{April 2026}

\begin{document}
\maketitle

\tableofcontents
\newpage

% ======================================================================
\section{Introduction}
\label{sec:intro}
% ======================================================================

This document describes the implementation of the Linear Parameter Varying (LPV) baseline
model for the ASMPT dual-gantry positioning stage, realised in the Linear Fractional
Representation (LFR) framework. The baseline serves as the non-trainable first-principles
component in the augmented model architecture proposed by Drenth~(2025, Ch.~5, eq.~5.1).

The baseline captures the Y-dependent inertia coupling of the gantry through the mass matrix
$M(Y)$, which is polynomial in the payload position~$Y$. Effects not present in the
Garc\'ia-Herreros first-principles equations---Coulomb friction, Coriolis/centripetal forces,
resonance dynamics---are excluded from the baseline by design (\decref{022}) and belong in the
augmentation component.

The document is structured as follows.
\Cref{sec:plant} defines the plant model and physical parameters.
\Cref{sec:lfr} derives the LFR formulation and the resolve-and-retain forward pass.
\Cref{sec:integration} explains the RK4 numerical integration choice.
\Cref{sec:framework} describes the interface with Jan's augmentation framework.
\Cref{sec:comparison} presents the comparison strategy for validation.
\Cref{sec:trajectory} discusses trajectory design.
\Cref{sec:matlab} describes the MATLAB data generation pipeline.
\Cref{sec:results} summarises the validation results.


% ======================================================================
\section{Plant Model}
\label{sec:plant}
% ======================================================================

\subsection{Equations of motion}

The dual-gantry is modelled as a 3-DOF system in logical coordinates
$q = [X,\, \Theta,\, Y]\T$, where $X$ is the rigid-body translation of the cross-arm,
$\Theta$ is the rotational deflection angle, and $Y$ is the payload position along the
cross-arm. The equation of motion is:
\begin{equation}
\label{eq:eom}
M(Y)\, \qddot = -K\, q - C\, \qdot + u,
\end{equation}
where $u = [F_X,\, \tau_\Theta,\, F_Y]\T$ are the generalised forces in logical coordinates.
The mass matrix $M(Y)$ depends on the scheduling variable $Y$ (payload position); the
stiffness matrix $K$ and viscous damping matrix $C$ are constant.

\subsection{Physical parameters}
\label{sec:params}

All parameters are taken from \texttt{kamtin-fp-model/03 Simulink gantry/main.m},
which is the immutable ground truth (\decref{002}). \Cref{tab:params} lists all values.

\begin{table}[ht]
\centering
\caption{Physical parameters of the dual-gantry model.}
\label{tab:params}
\begin{tabular}{llrl}
\toprule
Symbol & Description & Value & Unit \\
\midrule
$m_b$  & Mass of moving cross-arm        & 22.8   & kg \\
$m_h$  & Mass of payload (Y-axis)        & 10.1   & kg \\
$m_1$  & Mass of actuator X1             & 10.2   & kg \\
$m_2$  & Mass of actuator X2             & 10.7   & kg \\
$J_b$  & Rotary inertia of cross-arm     & 1.0    & kg\,m$^2$ \\
$J_h$  & Rotary inertia of payload       & 0.05   & kg\,m$^2$ \\
$c_{g1}$ & Viscous friction X1           & 14.5   & N/(m/s) \\
$c_{g2}$ & Viscous friction X2           & 20.3   & N/(m/s) \\
$c_y$    & Viscous friction Y            & 10.0   & N/(m/s) \\
$c_{b1}$ & Viscous friction joint 1      & 9.0    & Nm/(rad/s) \\
$c_{b2}$ & Viscous friction joint 2      & 9.0    & Nm/(rad/s) \\
$k_{b1}$ & Stiffness joint 1             & 1987.5 & Nm/rad \\
$k_{b2}$ & Stiffness joint 2             & 1987.5 & Nm/rad \\
$L_b$    & Length of cross-arm           & 0.725  & m \\
$d$      & Distance cross-arm to payload & 0.1    & m \\
\bottomrule
\end{tabular}
\end{table}

Python implementation: \pyfile{lpv\_lfr\_baseline/physics.py}, lines 27--46.

\subsection{Mass matrix decomposition}
\label{sec:mass_matrix}

The mass matrix is polynomial in $Y$:
\begin{equation}
\label{eq:MY}
M(Y) = M_0 + M_1 Y + M_2 Y^2,
\end{equation}
with the full $3 \times 3$ form:
\begin{equation}
\label{eq:MY_full}
M(Y) = \begin{bmatrix}
m_1 + m_2 + m_b + m_h & \frac{(m_1 - m_2) L_b}{2} - m_h Y & 0 \\[4pt]
\frac{(m_1 - m_2) L_b}{2} - m_h Y & J_b + J_h + \frac{(m_1 + m_2) L_b^2}{4} + m_h d^2 + m_h Y^2 & -m_h d \\[4pt]
0 & -m_h d & m_h
\end{bmatrix}.
\end{equation}
The decomposition by polynomial degree in $Y$:
\begin{align}
\label{eq:M0}
M_0 &= \begin{bmatrix}
m_1 + m_2 + m_b + m_h & \frac{(m_1 - m_2) L_b}{2} & 0 \\[3pt]
\frac{(m_1 - m_2) L_b}{2} & J_b + J_h + \frac{(m_1 + m_2) L_b^2}{4} + m_h d^2 & -m_h d \\[3pt]
0 & -m_h d & m_h
\end{bmatrix}, \\[6pt]
\label{eq:M1}
M_1 &= \begin{bmatrix}
0 & -m_h & 0 \\
-m_h & 0 & 0 \\
0 & 0 & 0
\end{bmatrix}, \\[6pt]
\label{eq:M2}
M_2 &= \begin{bmatrix}
0 & 0 & 0 \\
0 & m_h & 0 \\
0 & 0 & 0
\end{bmatrix}.
\end{align}
Python implementation: \pyfile{physics.py}, lines 62--78.

\subsection{Constant matrices}

The viscous damping matrix $C$ (no Y-dependence):
\begin{equation}
\label{eq:C}
C = \begin{bmatrix}
c_{g1} + c_{g2} & \frac{(c_{g1} - c_{g2}) L_b}{2} & 0 \\[3pt]
\frac{(c_{g1} - c_{g2}) L_b}{2} & c_{b1} + c_{b2} + \frac{(c_{g1} + c_{g2}) L_b^2}{4} & 0 \\[3pt]
0 & 0 & c_y
\end{bmatrix}.
\end{equation}
The stiffness matrix $K$ (no Y-dependence):
\begin{equation}
\label{eq:K}
K = \begin{bmatrix}
0 & 0 & 0 \\
0 & k_{b1} + k_{b2} & 0 \\
0 & 0 & 0
\end{bmatrix}.
\end{equation}
Note that $K$ has zero rows for X and Y---these are rigid-body modes with no spring restoring
force in those directions.

\subsection{What is included and excluded}

\paragraph{Included:}
\begin{itemize}[nosep]
\item Y-dependent inertia coupling through $M(Y)$
\item Viscous damping $C$ (constant)
\item Joint stiffness $K$ (constant)
\end{itemize}

\paragraph{Excluded (by design, \decref{022}):}
\begin{itemize}[nosep]
\item \textbf{Coulomb friction}: The Simulink Simscape model has Coulomb friction gains
  (\texttt{cc1}, \texttt{cc2}, \texttt{ccy}) present in the subsystem XML, but they carry
  \texttt{<P Name="Commented">on</P>}---they are \emph{disabled}. Verified by inspecting the
  \texttt{.slx} ZIP archive: \texttt{system\_47.xml} lists the Coulomb gain blocks as commented
  out. The quasi-LPV model \texttt{gantrySystem.m} (which produces the \texttt{q1} reference)
  does not include Coulomb friction.
\item \textbf{Coriolis/centripetal forces}: Not present in \texttt{gantrySystem.m}. These
  forces arise from the full nonlinear multibody equations (Simscape includes them) but are
  absent from the linearised quasi-LPV ODE. They belong in the augmentation (\decref{022}).
\item \textbf{Resonance dynamics}: Higher-order flexible dynamics not captured by the 3-DOF
  rigid-body model. To be addressed in the augmentation (\decref{024}).
\end{itemize}

\subsection{Coordinate systems}
\label{sec:coords}

The physics equations~\eqref{eq:eom} are formulated in \emph{logical coordinates}
$q = [X, \Theta, Y]\T$. However, experimental data from the real gantry is measured in
\emph{stage coordinates} $q_s = [X_1, X_2, Y]\T$ (encoder positions and amplifier forces).
The model operates in stage coordinates at the external interface (\decref{006}).

The transform between coordinate systems is given by the matrix $P$:
\begin{equation}
\label{eq:P}
P = \begin{bmatrix}
1 & 1 & 0 \\
L_b/2 & -L_b/2 & 0 \\
0 & 0 & 1
\end{bmatrix}
= \begin{bmatrix}
1 & 1 & 0 \\
0.3625 & -0.3625 & 0 \\
0 & 0 & 1
\end{bmatrix}.
\end{equation}
The relationships are:
\begin{align}
u_{\text{logical}} &= P \cdot u_{\text{stage}}, \label{eq:u_transform} \\
y_{\text{stage}} &= P\T \cdot y_{\text{logical}}. \label{eq:y_transform}
\end{align}
In the batched row-vector convention used by the Python code:
\begin{align}
\texttt{u\_logical} &= \texttt{u\_stage} \;\texttt{@}\; P\T, \\
\texttt{y\_stage} &= \texttt{y\_logical} \;\texttt{@}\; P.
\end{align}
Python implementation: \pyfile{physics.py}, lines 99--108.

\subsection{M(Y) invertibility}
\label{sec:invertibility}

The forward pass requires $M(Y)\inv$ at every timestep. $M(Y)$ is a physical mass matrix and
therefore symmetric positive definite for all physically realisable configurations. This was
verified numerically (\decref{021}): singular values of $M(Y)$ remain bounded away from zero
for $Y \in [-0.4,\, 0.4]\;\text{m}$ (the full operational range with margin). The determinant
$\det(M(Y))$ was compared against MATLAB at 50 $Y$ values, matching to better than
$10^{-8}$ (\pyfile{physics.py}, Check~2).


% ======================================================================
\section{LFR Formulation}
\label{sec:lfr}
% ======================================================================

\subsection{Why LFR?}

The supervisor confirmed (\decref{005}, updated 2026-03-22) that the baseline should use the
LFR structure. Drenth (2025, Ch.~5, eq.~5.1) assumes the baseline is available in LPV-LFR
form. The LFR parameterisation allows the learned augmentation correction to vary with $Y$ in
a principled way through the $\Delta$ block, and it separates the LTI dynamics from the
parameter-varying scheduling structure.

\subsection{The scheduling dependence}

From~\eqref{eq:eom}, the continuous-time state-space form is:
\begin{equation}
\label{eq:xdot}
\xdot = \begin{bmatrix} \qdot \\ \qddot \end{bmatrix}
= \begin{bmatrix} \qdot \\ M(Y)\inv\bigl(-K q - C \qdot + u\bigr) \end{bmatrix}
= \underbrace{\begin{bmatrix} 0 & I \\ -M(Y)\inv K & -M(Y)\inv C \end{bmatrix}}_{A_c(Y)}
x + \underbrace{\begin{bmatrix} 0 \\ M(Y)\inv \end{bmatrix}}_{B_c(Y)} u.
\end{equation}
Because $M(Y) = M_0 + M_1 Y + M_2 Y^2$ is polynomial in $Y$, the inverse $M(Y)\inv$ is
\emph{rational} in $Y$---not polynomial, not affine. This rational dependence is exactly what
the LFR is designed to handle: it factors the rational function into an LTI part (constant
matrices) and a structured uncertainty/scheduling block $\Delta(Y)$.

\subsection{The resolve-and-retain forward pass}
\label{sec:forward_pass}

The forward pass is implemented in \pyfile{lfr\_forward.py} as a 7-step sequence. The key
insight is the \emph{resolve-and-retain} approach: the algebraic loop created by $M(Y)\inv$
is resolved by solving the linear system, and the intermediate latent variables $z$ and $w$
are retained as explicit tensors (not discarded). This is what distinguishes the LFR form from
a simple collapsed evaluation of $A_c(Y) x + B_c(Y) u$---the latent signals are the
connection points where the augmentation network will attach.

\paragraph{Step 1: Build $M(Y)$ per batch.}
\begin{equation}
M(Y) = M_0 + M_1 Y + M_2 Y^2 \quad \in \R^{B \times 3 \times 3}
\end{equation}
where $B$ is the batch size.

\paragraph{Step 2: Compute net generalised force.}
\begin{equation}
f_{\text{net}} = -K q - C \qdot + u \quad \in \R^{B \times 3}
\end{equation}
where $q = x_{0:3}$ and $\qdot = x_{3:6}$ are the position and velocity components of the
state vector $x \in \R^6$.

\paragraph{Step 3: Resolve the algebraic loop.}
\begin{equation}
\label{eq:v_solve}
v = M(Y)\inv f_{\text{net}} \quad \in \R^{B \times 3}
\end{equation}
Implemented via \texttt{torch.linalg.solve(M\_Y, fnet)}, which is numerically stable and
differentiable. This is the acceleration $\qddot = M(Y)\inv f_{\text{net}}$.

The solve is the resolution of the algebraic loop: $M(Y)$ depends on $Y$, which is a state,
and $v$ feeds back into the state derivative. Solving the linear system computes the
unique $v$ that satisfies $M(Y) v = f_{\text{net}}$ without any iterative procedure---the
system is always well-posed because $M(Y)$ is invertible (\cref{sec:invertibility}).

\paragraph{Step 4: Latent variable chain from $\Delta(Y) = Y \cdot I_6$.}
\begin{align}
v_1 &= Y \cdot v \quad \in \R^{B \times 3}, \\
v_2 &= Y \cdot v_1 = Y^2 \cdot v \quad \in \R^{B \times 3}.
\end{align}
These are the scheduling-dependent quantities that carry the $Y$ and $Y^2$ dependence of
$M(Y)\inv$ into the LFR structure.

\paragraph{Step 5: Assemble LFR latent signals.}
\begin{align}
z &= \begin{bmatrix} v \\ v_1 \end{bmatrix} \in \R^{B \times 6}, \label{eq:z_def} \\
w &= \begin{bmatrix} v_1 \\ v_2 \end{bmatrix} \in \R^{B \times 6}. \label{eq:w_def}
\end{align}
These satisfy the fundamental LFR relationship:
\begin{equation}
\label{eq:delta}
\boxed{w = \Delta(Y) \cdot z, \qquad \Delta(Y) = Y \cdot I_6.}
\end{equation}
This is verified exactly (entry-wise, zero error) in \pyfile{lfr\_forward.py}, Check~5.
The $\Delta(Y) = Y \cdot I_6$ structure means the scheduling block is a simple scalar
multiplication broadcast across all 6 latent channels.

\paragraph{Why $z$ and $w$ are retained.}
The augmentation network connects \emph{between} $z$ and $w$. In the augmented model
(Drenth Ch.~5, eq.~5.1--5.2), the augmentation $\Delta^a$ block receives $z$ and produces a
correction to $w$. If we collapsed $z$ and $w$ into the state derivative (as in a simple
$A_c(Y) x + B_c(Y) u$ evaluation), these connection points would not exist. Retaining them as
explicit tensors is the defining feature of the LFR implementation.

\paragraph{Step 6: State derivative---direct from physics (\decref{026}).}
\begin{equation}
\label{eq:xdot_direct}
\xdot = \begin{bmatrix} x_{3:6} \\ v \end{bmatrix} = \begin{bmatrix} \qdot \\ \qddot \end{bmatrix} \in \R^{B \times 6}.
\end{equation}
This uses the physical identity directly: the first 3 state derivatives are the velocities
(already in the state), and the last 3 are the accelerations $v = M(Y)\inv f_{\text{net}}$
(computed in Step~3).

\textbf{Why not the G-matrix form?} An earlier implementation used $\xdot = G_{Ax} x +
G_{Bw} w + G_{Bu} u$, where $G_{Ax}$, $G_{Bw}$, $G_{Bu}$ are constant LFR matrices.
This was removed (\decref{026}) because the G-matrices are precomputed using $M_0\inv$ and
become stale if physical parameters are updated during training (parameter estimation). The
direct physics expression~\eqref{eq:xdot_direct} is always correct for whatever $M_0, M_1,
M_2, K, C$ are passed to the forward function at that call. The two forms are algebraically
identical---verified in \pyfile{lfr\_forward.py}, Check~2---but the direct form eliminates
the silent-staleness risk.

\paragraph{Step 7: Output.}
\begin{equation}
y = x_{0:3} = q \quad \in \R^{B \times 3} \quad \text{(logical positions)}.
\end{equation}
The output is the position part of the state. The conversion to stage coordinates
$y_{\text{stage}} = P\T y$ is applied by the caller (\pyfile{lfr\_simulate.py}).

\subsection{Self-scheduling: quasi-LPV}
\label{sec:self_scheduling}

The scheduling variable $Y$ is not an exogenous signal---it is extracted from the state:
$Y = x_2$ (the third state component in logical coordinates, corresponding to the payload
position). This makes the system \emph{quasi-LPV}: the scheduling variable is a function of
the state, not an independent external input.

In the forward pass, $Y$ is extracted as \texttt{Y = x[:, 2]} before calling
\texttt{lfr\_forward}. During RK4 integration (\cref{sec:integration}), $Y$ is re-extracted
at each of the 4 sub-step evaluations from the intermediate state, ensuring the scheduling
tracks the evolving state throughout the integration step.

\subsection{Mathematical verification}
\label{sec:lfr_verification}

Five checks are implemented in \pyfile{lfr\_forward.py} (\texttt{\_\_main\_\_}):
\begin{enumerate}[nosep]
\item \textbf{Loop resolution residual}: $\|M(Y) v - f_{\text{net}}\|_\infty < 10^{-12}$
  for 5 test $Y$ values.
\item \textbf{CT vector field}: $\xdot$ matches the collapsed $A_c(Y) x + B_c(Y) u$ to
  $< 10^{-12}$ for 5 $Y$ values.
\item \textbf{Autograd through $Y$}: gradient $\partial L / \partial Y$ exists after backward
  pass through \texttt{torch.linalg.solve}.
\item \textbf{Autograd through $x$}: gradient $\partial L / \partial x$ exists after backward
  pass.
\item \textbf{LFR signal structure}: $w = Y \cdot z$ verified entry-wise with zero error for
  all 5 $Y$ values.
\end{enumerate}
All checks pass, confirming the forward pass is mathematically correct and fully
differentiable.


% ======================================================================
\section{Numerical Integration}
\label{sec:integration}
% ======================================================================

\subsection{RK4 choice (\decref{018})}

The supervisor directed: ``use RK4, not Euler discretization. Better to not precompute.''
and ``write up the CT conversion. Don't do discretisation first, will get messy.'' The model
is therefore maintained in continuous time and integrated using a classical 4th-order
Runge--Kutta (RK4) method with fixed step size $h = \ts = 1/\fs$.

\subsection{RK4 formula with self-scheduling}

Given state $x[k]$ and input $u[k]$ (held constant over the step), one RK4 step computes:
\begin{align}
k_1 &= f\bigl(x[k],\; u[k],\; Y = x[k]_2\bigr), \label{eq:rk4_k1} \\
k_2 &= f\bigl(x[k] + \tfrac{h}{2} k_1,\; u[k],\; Y = (x[k] + \tfrac{h}{2} k_1)_2\bigr), \label{eq:rk4_k2} \\
k_3 &= f\bigl(x[k] + \tfrac{h}{2} k_2,\; u[k],\; Y = (x[k] + \tfrac{h}{2} k_2)_2\bigr), \label{eq:rk4_k3} \\
k_4 &= f\bigl(x[k] + h\, k_3,\; u[k],\; Y = (x[k] + h\, k_3)_2\bigr), \label{eq:rk4_k4} \\
x[k\!+\!1] &= x[k] + \frac{h}{6}\bigl(k_1 + 2 k_2 + 2 k_3 + k_4\bigr), \label{eq:rk4_update}
\end{align}
where $f(\cdot)$ is the \texttt{lfr\_forward} function returning $\xdot$, and the subscript
$(\cdot)_2$ denotes extraction of the 3rd state component (Y-position in logical coordinates).

The critical detail is that $Y$ is \emph{re-extracted from the intermediate state} at each
sub-step evaluation ($k_1$ through $k_4$). This is the self-scheduling aspect: the scheduling
variable tracks the state as it evolves within the integration step, rather than being frozen
at $Y[k]$.

Python implementation: \pyfile{lfr\_simulate.py}, lines 66--100.

The $z$ and $w$ latent signals are recorded from the $k_1$ evaluation only (at the start of
the step, $x[k]$). This is consistent with the LFR interpretation: the latent signals
represent the scheduling structure at the current state, and the RK4 sub-steps are internal
to the integrator.

\subsection{Why not ZOH pre-discretisation?}

The zero-order hold (ZOH) approach pre-computes discrete matrices $A_d(Y)$, $B_d(Y)$ via
matrix exponential and then applies $x[k+1] = A_d x[k] + B_d u[k]$. This was used for
Steps~1--2 validation (\decref{012}) where exact MATLAB matrix comparison was the goal. For
the training loop, the supervisor explicitly directed against pre-discretisation (\decref{018}).

ZOH holds the input constant within each interval but says nothing about how the ODE is
integrated inside the interval. RK4 is the integration method used inside that interval. The
two are not alternatives to the same problem---they address different aspects of discretisation.

\subsection{Why not Euler?}

Euler (forward) is a 1st-order method with local truncation error $O(h^2)$ per step and
global error $O(h)$. RK4 has local truncation error $O(h^5)$ per step and global error
$O(h^4)$. At $h = 50\;\mu\text{s}$ ($\fs = 20\;\text{kHz}$), this is a dramatic accuracy
difference. The supervisor confirmed: ``use RK4 not Euler.''

\subsection{Why not ODE45?}

ODE45 uses adaptive step sizes (Dormand--Prince), which is incompatible with a fixed sampling
period in a discrete control loop. The \texttt{ode4} variant forces a fixed step, but that is
equivalent to RK4 directly.

\subsection{Step size and sample rate}

The step size is $h = \ts = 1/\fs = 50\;\mu\text{s}$, corresponding to $\fs = 20\;\text{kHz}$.
This matches the Position Loop Time Interval (PLTI) from the AccurET hardware specification
(\texttt{AccurET-Oper\&Soft-VerV.pdf}). The PLTI is the rate at which the position control
loop executes on the real hardware, making it the natural sampling period for plant modelling.

Note: the original \texttt{main.m} used $\fs = 16\;\text{kHz}$, which was a discrepancy with
the hardware specification. All files have been updated to $20\;\text{kHz}$ (\decref{012},
resolved note).

\subsection{Numerical precision}

All physics computations use \texttt{float64} (double precision). This is required to avoid
error accumulation over long trajectories---a typical simulation of $1.75\;\text{s}$ at
$20\;\text{kHz}$ involves $35{,}000$ integration steps. With \texttt{float32}, the accumulated
rounding error would degrade the trajectory after several thousand steps.


% ======================================================================
\section{Jan's Augmentation Framework Interface}
\label{sec:framework}
% ======================================================================

\subsection{Block base class contract}

Jan's augmentation framework (\texttt{model\_augmentation/fit\_systems/blocks.py}) defines a
\texttt{Block} base class with the interface:
\begin{lstlisting}[language=Python]
class Block(nn.Module):
    def __init__(self, nz, nw):
        self.nz = nz   # input dimension
        self.nw = nw   # output dimension

    def forward(self, z: Tensor) -> Tensor:
        # z shape: (batch, nz, 1)
        # returns: (batch, nw, 1)
\end{lstlisting}
Every block receives a 3D tensor $z \in \R^{B \times n_z \times 1}$ and returns
$w \in \R^{B \times n_w \times 1}$. The trailing dimension of~1 is a framework convention
(column-vector semantics in the batch dimension).

\subsection{LFRBaselineBlock: nz = 9, nw = 18}

The baseline block (\pyfile{lfr\_block.py}) implements the \texttt{Block} interface with:
\begin{align}
n_z &= 9 \quad (n_x = 6 \text{ state} + n_u = 3 \text{ stage-coord input}), \\
n_w &= 18 \quad (x_{\text{next}} = 6 + z_{\text{lfr}} = 6 + w_{\text{lfr}} = 6).
\end{align}

\paragraph{Input packing:}
\begin{equation}
z_{\text{in}} = \begin{bmatrix} x \\ u_{\text{stage}} \end{bmatrix} \in \R^{B \times 9 \times 1},
\qquad
z_{\text{in}}[\,:\,,\, 0\!:\!6\,,\, :] = x, \quad
z_{\text{in}}[\,:\,,\, 6\!:\!9\,,\, :] = u_{\text{stage}}.
\end{equation}

\paragraph{Output packing:}
\begin{equation}
w_{\text{out}} = \begin{bmatrix} x_{\text{next}} \\ z_{\text{lfr}} \\ w_{\text{lfr}} \end{bmatrix} \in \R^{B \times 18 \times 1},
\quad
\begin{cases}
w_{\text{out}}[\,:\,,\, 0\!:\!6\,,\, :] = x_{\text{next}} & \text{(next state, logical coords)}, \\
w_{\text{out}}[\,:\,,\, 6\!:\!12\,,\, :] = z_{\text{lfr}} & \text{(LFR latent $z$)}, \\
w_{\text{out}}[\,:\,,\, 12\!:\!18\,,\, :] = w_{\text{lfr}} & \text{(LFR latent $w$)}.
\end{cases}
\end{equation}

\subsection{Stateless design}

The block is \emph{stateless}: it does not store or update any internal state variable.
State management is handled externally by Jan's \texttt{Interconnect} class, which routes
$x_{\text{next}}$ from the block output back to the block input at each timestep. This is
critical: storing \texttt{self.x} inside the block and updating it in \texttt{forward()} would
break the \texttt{Interconnect}'s state management assumptions.

\subsection{float32/float64 boundary}

Jan's framework operates in \texttt{float32} (standard for neural network training). The
physics computation requires \texttt{float64} for numerical accuracy (\cref{sec:integration}).
The boundary is managed by two cast operations in \texttt{forward()}:
\begin{lstlisting}[language=Python]
# Entry: float32 -> float64
z_f64 = z_in.squeeze(-1).double()

# [... all physics in float64 ...]

# Exit: float64 -> float32
return w_f64.float().unsqueeze(-1)
\end{lstlisting}
To run the entire pipeline in \texttt{float64} (e.g., for validation), remove these two cast
lines---no other file needs changing.

\subsection{Naming convention}

Jan's API uses `$z$' for block input and `$w$' for block output. Our LFR also has latent
signals called $z$ and $w$. To avoid confusion, the LFR latent signals are called
$z_{\text{lfr}}$ and $w_{\text{lfr}}$ inside the block code.

\subsection{How augmentation connects}

In the augmented model, the augmentation $\Delta^a$ network sits between the $w_{\text{lfr}}$
and $z_{\text{lfr}}$ slots in the \texttt{Interconnect} wiring. The \texttt{Interconnect}
reads $z_{\text{lfr}}$ and $w_{\text{lfr}}$ from the baseline block's output, routes them
through the augmentation network, and combines the result with the baseline's state update.

\subsection{Coordinate transform at the boundary (\decref{027})}

The block output $y = x_{0:3}$ is in logical coordinates $[X, \Theta, Y]$, but the training
data and reference signals are in stage coordinates $[X_1, X_2, Y]$. The coordinate transform
$P\T$ must be applied somewhere. The correct place is the \texttt{Interconnect}'s output
connection matrix $S_y$, not inside the block (which would change the $n_w = 18$ contract):
\begin{equation}
S_y = P\T \cdot S_{\text{select}}(0\!:\!3,\; 18),
\end{equation}
where $S_{\text{select}}$ is a $(3 \times 18)$ selection matrix that picks the first 3
components of $w_{\text{out}}$, and $P\T$ transforms from logical to stage coordinates.

\subsection{Physical parameters as attributes}

Physical parameters ($M_0$, $M_1$, $M_2$, $K$, $C$, $P$, $\ts$) are stored as plain
attributes. For GPU training, these should be replaced with
\texttt{self.register\_buffer(name, tensor)} so tensors move automatically with the model
to the correct device. No other change is needed.


% ======================================================================
\section{Comparison Strategy}
\label{sec:comparison}
% ======================================================================

\subsection{Open-loop replay rationale}

All models are compared by replaying the \emph{same} force sequence $u_{q1}[k]$ in open loop.
This isolates model quality from controller compensation. In a closed-loop comparison, a better
controller could mask a worse model (and vice versa), confounding the assessment.

\subsection{Reference signal: why q1}
\label{sec:why_q1}

Three MATLAB outputs are available from the Simulink simulation:

\begin{description}[nosep]
\item[\texttt{q1}] CT quasi-LPV, integrated by Simulink using \texttt{gantrySystem.m}.
  Self-schedules $Y = x(3)$ at every ODE step. $M(Y)$ varies continuously. No Coulomb, no
  Coriolis. \textbf{Primary comparison target.}
\item[\texttt{q}] Simscape full nonlinear model. Coulomb friction disabled (verified by SLX
  XML inspection). Includes Coriolis/centripetal. \textbf{Secondary comparison target.}
\item[\texttt{q3}] Frozen LTI \texttt{lsim} at $Y = 0.3\;\text{m}$. Computed post-hoc in
  \texttt{main.m} via \texttt{lsim(T, ...)}. Uses its \emph{own} forces $u_{q3}$.
\end{description}

\texttt{q1} is chosen as the primary reference because it integrates the \emph{same} CT
quasi-LPV ODE as the Python implementation---same $M(Y)$, same $C$, same $K$, no Coulomb,
no Coriolis. The Python LPV model driven by $u_{q1}$ should match \texttt{q1} to numerical
precision (integration method mismatch only: RK4 fixed-step vs MATLAB \texttt{ode45} adaptive).

\paragraph{Why NOT q3.}
\texttt{q3} runs in its own closed loop: $u_{q3} = C_{fb}(r - q3)$. These forces are
different from $u_{q1}$ because the frozen LTI tracks the reference differently. Comparing
the Python frozen model driven by $u_{q1}$ against \texttt{q3} driven by $u_{q3}$ confounds
the model quality assessment with the force mismatch.

\paragraph{Why NOT Simscape as primary.}
The Simscape model includes Coriolis/centripetal forces that our baseline intentionally
excludes (\decref{022}). It also runs with its own forces $u_q = C_{fb}(r - q)$, different
from $u_{q1}$. Using Simscape as the primary reference would mix the ``Coriolis gap'' with
the ``force mismatch,'' making it impossible to isolate the LPV scheduling benefit.

\subsection{Frozen LTI baseline}

The frozen LTI comparison is implemented in Python (\texttt{simulate\_frozen} in
\pyfile{validate\_lfr.py}): the same simulation code as the LPV model, but with $Y$ held
constant at $Y = 0.3\;\text{m}$ (the design operating point from \texttt{main.m}). This is
\emph{not} the same as MATLAB's \texttt{q3}---it uses the same forces $u_{q1}$ as the LPV
model, isolating the scheduling benefit.

\subsection{Simscape secondary comparison}

The Simscape model is compared as a secondary reference using its own matched forces $u_q$.
The Python LPV model is re-simulated with $u_q$ as input (not $u_{q1}$), ensuring a fair
comparison. The gap between Python LPV and Simscape (both driven by $u_q$) quantifies the
Coriolis effect---physics that our baseline intentionally does not model and that the
augmentation must learn.

\subsection{Metrics}

\begin{description}[nosep]
\item[BFR (Best Fit Rate)] The field-standard, scale-invariant quality metric:
\begin{equation}
\label{eq:bfr}
\text{BFR} = \max\!\left(1 - \frac{\|y - \hat{y}\|_2}{\|y - \bar{y}\|_2},\; 0\right) \times 100\%,
\end{equation}
where $y$ is the reference, $\hat{y}$ is the model prediction, and $\bar{y}$ is the
reference mean. BFR = 100\% means perfect prediction; BFR = 0\% means the model is no
better than predicting the mean.
\item[$\max|e|$] Maximum absolute error---engineering bound on worst-case deviation.
\item[RMSE] Root mean square error---overall deviation magnitude.
\item[$\text{Mean}|e|$] Mean absolute error---average deviation for completeness.
\end{description}


% ======================================================================
\section{Trajectory Design}
\label{sec:trajectory}
% ======================================================================

\subsection{Y motion: 0.3 to $-$0.3 m}

The test trajectory moves $Y$ from $+0.3\;\text{m}$ to $-0.3\;\text{m}$---a displacement of
$0.6\;\text{m}$ in the negative direction. This trajectory was chosen to maximise the frozen
LTI modelling error:

\begin{itemize}[nosep]
\item The $M_{01}$ entry (mass matrix off-diagonal coupling) is:
  $M_{01}(Y) = \frac{(m_1 - m_2) L_b}{2} - m_h Y$.
  This changes sign at $Y \approx 0.018\;\text{m}$. The trajectory crosses this sign-flip
  point, making the frozen model (evaluated at $Y = 0.3\;\text{m}$) qualitatively wrong for
  a significant portion of the motion.
\item At $Y = +0.3\;\text{m}$: $M_{01} = -3.21\;\text{kg\,m}$.
  At $Y = -0.3\;\text{m}$: $M_{01} = +3.17\;\text{kg\,m}$.
  The sign flip means the frozen model predicts the wrong direction of X--$\Theta$ coupling.
\end{itemize}

\subsection{X1, X2 at rest}

$X_1$ and $X_2$ references are held at zero throughout the trajectory. This decouples the
Coriolis effect (which is proportional to $\dot{X} \cdot \dot{Y}$): when $\dot{X} \approx 0$,
Coriolis forces vanish. This allows isolating the Y-varying inertia effect from
Coriolis/centripetal coupling.

\subsection{Hardware limits}

Motion parameters are verified against the ETEL TELICA datasheet
(\texttt{telica-xyz-0750-0800-data.pdf}):

\begin{table}[ht]
\centering
\caption{Motion profile parameters vs hardware limits.}
\label{tab:motion}
\begin{tabular}{lrrl}
\toprule
Parameter & Used & Hardware max & Fraction \\
\midrule
$v_{\max}$ & 1.0 m/s & 2.0 m/s & 50\% \\
$a_{\max}$ & 10 m/s$^2$ & 50 m/s$^2$ & 20\% \\
Stroke     & $\pm 300$ mm & $\pm 400$ mm & 75\% \\
\bottomrule
\end{tabular}
\end{table}

The motion profile uses a 3rd-order jerk-limited setpoint
(\texttt{thirdOrderSetpointETEL})---the same profile generator used by the real machine's
motion controller.

\subsection{Why this trajectory maximises the frozen LTI error}

The key is the $M_{01}$ sign flip at $Y \approx 0.018\;\text{m}$. The frozen LTI at
$Y = 0.3\;\text{m}$ has a negative $M_{01}$. As $Y$ moves past $0.018\;\text{m}$ toward
$-0.3\;\text{m}$, the true $M_{01}$ becomes positive---the frozen model predicts the wrong
sign for the inertia coupling. This creates a growing discrepancy that cannot be corrected by
simply scaling the frozen matrices. The LPV model, which re-evaluates $M(Y)$ at every step,
tracks this sign change correctly.


% ======================================================================
\section{MATLAB Data Generation Pipeline}
\label{sec:matlab}
% ======================================================================

All reference data used for Python validation is generated by three MATLAB scripts in
\texttt{Matlab-scripts/}. The scripts call functions from the immutable
\texttt{kamtin-fp-model/} directory but do not modify any file in it.
Run each script from the repository root with MATLAB's working directory set there.

\subsection{export\_lpv\_matrices.m --- discrete A(Y), B(Y) at 50 operating points}
\label{sec:export_matrices}

\textbf{Purpose}: Provides a dense grid of MATLAB-computed discrete-time state matrices for
Python matrix comparison (Task~2.4, \decref{016}).

\textbf{What it does}:
\begin{enumerate}[nosep]
\item Replicates the physics setup from \texttt{main.m} (lines 12--64): physical parameters,
  $C$, $K$, $P$ (identical values to \pyfile{physics.py}).
\item Sweeps $Y$ over 50 values: \texttt{linspace(-0.35, 0.35, 50)} m---the full operational
  range with margin.
\item At each $Y$: builds $M(Y)$, calls \texttt{getss(n, M, C, K)} to get the logical-coordinate
  continuous-time system, applies the $P$-transform to stage coordinates, then discretises with
  \texttt{c2d(sys\_stage, ts, 'zoh')}.
\item Stores $A(Y)$, $B(Y)$, $C$, $D$ per operating point, plus $\det(M(Y))$ as a
  physics health check.
\end{enumerate}

\textbf{Output file}: \texttt{Matlab-output/lpv\_matrices.mat}

\begin{table}[ht]
\centering
\caption{Variables in \texttt{lpv\_matrices.mat}.}
\label{tab:lpv_matrices}
\begin{tabular}{lll}
\toprule
Variable & Shape & Description \\
\midrule
\texttt{A\_all}   & $(6,6,50)$ & Discrete-time state matrix $A(Y)$ at each operating point \\
\texttt{B\_all}   & $(6,3,50)$ & Discrete-time input matrix $B(Y)$ at each operating point \\
\texttt{C\_all}   & $(3,6,50)$ & Output matrix (constant, stored per Y for convenience) \\
\texttt{D\_all}   & $(3,3,50)$ & Feedthrough matrix (identically zero) \\
\texttt{Y\_values} & $(50,1)$  & Y sweep values [m] \\
\texttt{det\_M}   & $(50,1)$  & $\det(M(Y))$ at each Y (physics health check) \\
\texttt{fs}, \texttt{ts} & scalar & Sample frequency and period \\
\bottomrule
\end{tabular}
\end{table}

\textbf{Used by}: \pyfile{validate\_lfr.py} Section~3 (Y-varying Bode), and
\pyfile{physics.py} Check~2 ($\det M$ comparison).

\textbf{Why 50 points instead of 5}: $M(Y)\inv$ is rational in $Y$ and could have
non-monotone error behaviour between sparse samples. A dense 50-point sweep allows plotting
error vs $Y$ to confirm uniformity across the full operational range.

\subsection{export\_lpv\_sim.m --- closed-loop varying-Y Simulink simulation}
\label{sec:export_sim}

\textbf{Purpose}: Provides the varying-$Y$ closed-loop trajectory that proves the LPV model
captures Y-dependent dynamics better than a frozen LTI.

\textbf{What it does}:
\begin{enumerate}[nosep]
\item Sets up identical physics to \texttt{main.m}.
\item Constructs the test trajectory: $Y$ moves from $+0.3$ to $-0.3\;\text{m}$ using a
  3rd-order jerk-limited profile (\texttt{thirdOrderSetpointETEL}); $X_1 = X_2 = 0$
  throughout.
\item Builds the discrete feedback controller $C_{fb}$ (rule-of-thumb, 100~Hz bandwidth) and
  the frozen LTI $G$ at $Y = 0.3\;\text{m}$ (needed by Simulink at runtime for the LTI
  block).
\item Runs the Simulink model \texttt{gantry\_2025a} via \texttt{sim(mdl, t(end))}.
  Three ToWorkspace blocks produce \texttt{q} (Simscape), \texttt{q1} (CT quasi-LPV), and
  \texttt{q2} (CT quasi-LPV + Coriolis; not exported).
\item Reconstructs the forces applied to each path post-hoc using \texttt{lsim}:
  $u_{q1} = C_{fb}(r - q_1)$ and $u_q = C_{fb}(r - q)$.
\item Saves all variables to \texttt{Matlab-output/lpv\_sim\_varying\_y.mat}.
\end{enumerate}

\textbf{Output file}: \texttt{Matlab-output/lpv\_sim\_varying\_y.mat}

\begin{table}[ht]
\centering
\caption{Variables in \texttt{lpv\_sim\_varying\_y.mat}.}
\label{tab:lpv_sim}
\begin{tabular}{lll}
\toprule
Variable & Shape & Description \\
\midrule
\texttt{t\_sim}        & $(N,1)$   & Time vector [s] \\
\texttt{fs}            & scalar    & Sample frequency = 20\,000 Hz \\
\texttt{r\_sim}        & $(N,3)$   & Reference $[X_{1,\text{ref}},\, X_{2,\text{ref}},\, Y_{\text{ref}}]$ [m] \\
\texttt{u\_q1}         & $(N,3)$   & Forces applied to q1 path: $C_{fb}(r - q_1)$ [N] \\
\texttt{u\_q}          & $(N,3)$   & Forces applied to Simscape path: $C_{fb}(r - q)$ [N] \\
\texttt{q1}            & $(N,3)$   & CT quasi-LPV output $[X_1, X_2, Y]$ [m] --- primary reference \\
\texttt{q\_simscape}   & $(N,3)$   & Simscape nonlinear output $[X_1, X_2, Y]$ [m] \\
\texttt{Y\_trajectory} & $(N,1)$   & $Y(t) = q_1[:,2]$ -- scheduling variable trajectory [m] \\
\bottomrule
\end{tabular}
\end{table}

\textbf{Why u\_q and u\_q1 are different}: The two paths run in their own closed loops
during the Simulink simulation. \texttt{q1} tracks the reference with forces $u_{q1} =
C_{fb}(r - q_1)$; the Simscape model tracks with forces $u_q = C_{fb}(r - q)$. Because the
two outputs differ slightly (quasi-LPV ODE vs full nonlinear Simscape), the forces differ.
For a fair Python--MATLAB comparison, the Python model must be driven by the \emph{same}
forces as the MATLAB reference it is compared against (\cref{sec:why_q1}).

\textbf{Why the save path uses fileparts}: The save path is resolved as
\texttt{fullfile(fileparts(mfilename('fullpath')), '..', 'Matlab-output')}, so it saves to
the correct directory regardless of MATLAB's current working directory. This was a bug fix:
the original \texttt{fullfile(pwd, 'Matlab-output')} saved to wherever MATLAB's \texttt{pwd}
happened to be, often in the wrong place.

\textbf{Coulomb friction verification}: The Simscape Coulomb friction gains (\texttt{cc1},
\texttt{cc2}, \texttt{ccy}) are defined as workspace variables in this script but are
disabled in the Simulink model. Verified by inspecting \texttt{gantry\_2025a.slx} as a ZIP
archive: the subsystem XML (\texttt{system\_47.xml}) lists the Coulomb gain blocks with
\texttt{<P Name="Commented">on</P>}, confirming they are commented out and have no effect.

\subsection{verify\_exports.m --- automated validation of both exports}
\label{sec:verify_exports}

\textbf{Purpose}: A single entry-point that runs both export scripts and checks their outputs
against a set of physical sanity checks. Run this after any change to either export script.

\textbf{What it checks}:

\paragraph{lpv\_matrices checks (6 checks):}
\begin{itemize}[nosep]
\item Array dimensions: $A$ is $(6,6,50)$, $B$ is $(6,3,50)$, $C$ is $(3,6,50)$, $D$ is $(3,3,50)$.
\item $Y$ values within $[-0.35, 0.35]\;\text{m}$.
\item $D$ matrix identically zero (max abs $< 10^{-12}$).
\item $\det(M(Y)) > 0$ at all 50 operating points (mass matrix positive definite).
\item $C$ matrix constant across all $Y$ values (max std $< 10^{-10}$).
\item All eigenvalues of $A(Y)$ strictly inside the unit circle (stable discrete system).
\end{itemize}

\paragraph{lpv\_sim\_varying\_y checks (9 checks):}
\begin{itemize}[nosep]
\item Consistent array dimensions.
\item $Y$ trajectory within physical machine range $[-0.4, 0.4]\;\text{m}$.
\item Initial $Y = 0.3\;\text{m}$ (tolerance $\pm 2\;\text{mm}$).
\item Final $Y = 0.1\;\text{m}$ (tolerance $\pm 2\;\text{mm}$).
  \emph{Note: the current trajectory goes to $-0.3\;\text{m}$; this check catches an
  outdated trajectory if verify\_exports.m is not updated to match.}
\item $F_Y$ RMS $> 1\;\text{N}$ (confirms feedback was active).
\item $X_1$, $X_2$ max deviation $< 10\;\mu\text{m}$ (X reference held at zero).
\item $Y$ tracking error in final hold period $< 50\;\mu\text{m}$.
\item Forces within ETEL peak limits: $F_{X1}, F_{X2} < 2000\;\text{N}$;
  $F_Y < 1420\;\text{N}$.
\end{itemize}

\subsection{Other files in Matlab-output/}

\begin{description}[nosep]
\item[\texttt{gantry\_G\_matrices.mat}] Discrete $A$, $B$, $C$, $D$ at the single operating
  point $Y = 0.3\;\text{m}$. Used by \pyfile{validate\_lfr.py} Section~2 (Bode comparison).
\item[\texttt{gantry\_input.mat}] Force input sequence $u[k]$ for the initial frozen LTI
  simulation. Used by \pyfile{lfr\_simulate.py} Check~3.
\item[\texttt{gantry\_q3\_lsim.mat}] MATLAB \texttt{lsim} output for the frozen LTI at
  $Y = 0.3\;\text{m}$. Used by \pyfile{lfr\_simulate.py} Check~3.
\item[\texttt{gantry\_q\_simscape.mat}] Simscape output from an earlier single-point run (not
  the varying-$Y$ trajectory). Superseded by \texttt{lpv\_sim\_varying\_y.mat} for LPV
  validation purposes.
\item[\texttt{workspace\_snapshot.txt}] Text dump of key MATLAB workspace variables at the
  time of export. Useful for auditing that export scripts ran with the intended parameters.
\end{description}


% ======================================================================
\section{Results}
\label{sec:results}
% ======================================================================

\subsection{Summary table}

\Cref{tab:results} shows the validation results from all 5 sections of
\pyfile{validate\_lfr.py}. All models are driven by the same forces $u_{q1}$ (open-loop
replay).

\begin{table}[ht]
\centering
\caption{Validation results: Python LPV and frozen LTI vs MATLAB q1 reference.
  All models driven by $u_{q1}$ forces. Y range: $-0.3$ to $+0.3\;\text{m}$.}
\label{tab:results}
\begin{tabular}{llrrrr}
\toprule
Model & Channel & BFR [\%] & RMSE [m] & Max$|e|$ [m] & Mean$|e|$ [m] \\
\midrule
Python LPV    & X1 & 100.0000 & $1.18 \times 10^{-14}$ & $2.60 \times 10^{-14}$ & $7.66 \times 10^{-15}$ \\
Python LPV    & X2 & 100.0000 & $1.17 \times 10^{-14}$ & $2.57 \times 10^{-14}$ & $7.61 \times 10^{-15}$ \\
Python LPV    & Y  & 100.0000 & $8.64 \times 10^{-14}$ & $2.25 \times 10^{-13}$ & $5.39 \times 10^{-14}$ \\
\midrule
Python frozen & X1 & 52.84 & $1.69 \times 10^{-7}$ & $6.04 \times 10^{-7}$ & $1.02 \times 10^{-7}$ \\
Python frozen & X2 & 36.49 & $1.52 \times 10^{-7}$ & $5.10 \times 10^{-7}$ & $9.27 \times 10^{-8}$ \\
Python frozen & Y  & 100.00 & $2.62 \times 10^{-8}$ & $9.31 \times 10^{-8}$ & $1.58 \times 10^{-8}$ \\
\bottomrule
\end{tabular}
\end{table}

\subsection{Natural frequencies vs Y}

The eigenvalue analysis of $A_c(Y)$ across $Y \in [-0.4, +0.4]\;\text{m}$ reveals that only
the $\Theta$ mode (cross-arm rotation) is oscillatory. The X and Y modes are overdamped
(non-oscillatory)---their eigenvalues are purely real.

The oscillatory natural frequency varies from $4.41\;\text{Hz}$ (at $|Y| = 0.4\;\text{m}$) to
$5.13\;\text{Hz}$ (at $Y = 0\;\text{m}$). This is because the effective rotational inertia
$J_{\text{eff}}(Y) = J_b + J_h + (m_1 + m_2) L_b^2/4 + m_h d^2 + m_h Y^2$ increases with
$|Y|$ (the payload acts as a point mass at distance $Y$ from the rotation axis), lowering the
natural frequency $\omega_n \sim \sqrt{k / J_{\text{eff}}}$.

\subsection{Key findings}

\begin{enumerate}
\item \textbf{LPV model matches MATLAB to machine precision}: BFR = 100.0\% across all
  channels, with maximum error $2.25 \times 10^{-13}\;\text{m}$. This confirms the Python
  implementation correctly reproduces the MATLAB quasi-LPV ODE.

\item \textbf{Frozen LTI fails on X channels}: BFR drops to 52.8\% (X1) and 36.5\% (X2)
  when $Y$ varies from $+0.3$ to $-0.3\;\text{m}$. The Y channel remains at 100\% because
  $M_{22} = m_h$ is constant---Y dynamics do not depend on $Y$ in the mass matrix.

\item \textbf{Frozen LTI failure is asymmetric}: X2 (BFR = 36.5\%) degrades more than X1
  (BFR = 52.8\%). This is because the mass asymmetry $(m_1 \neq m_2)$ makes the off-diagonal
  coupling $M_{01}$ asymmetric: the frozen model's error in $M_{01}$ affects X1 and X2
  differently through the $P$-transform.

\item \textbf{LPV scheduling benefit is substantial}: The 47--64 percentage-point BFR
  improvement on X1/X2 demonstrates that tracking $M(Y)$ is essential for accurate prediction
  when $Y$ varies over its operational range.
\end{enumerate}


% ======================================================================
\section*{Decision Reference}
% ======================================================================

Key design decisions referenced in this document, logged in
\texttt{docs/decisions.md}:

\begin{description}[nosep,font=\ttfamily]
\item[D-002] MATLAB files in \texttt{kamtin-fp-model/} are immutable ground truth.
\item[D-005] LFR structure confirmed for the LPV augmentation (supervisor, 2026-03-20).
\item[D-006] Python implementation uses stage coordinates at the external interface.
\item[D-012] LPV discretisation: frozen ZOH for validation, CT+RK4 for training.
\item[D-018] CT model with RK4 integration at fixed step (supervisor directive).
\item[D-020] Method~2 (state-space form) chosen for rational LPV dependency.
\item[D-022] Non-baseline physics go in augmentation, not in baseline.
\item[D-026] Direct physics replaces G-matrix form in forward pass.
\item[D-027] $P\T$ coordinate transform embedded in Interconnect $S_y$ matrix.
\end{description}


% ======================================================================
\section*{Source File Reference}
% ======================================================================

\begin{description}[nosep,font=\ttfamily]
\item[physics.py] Physical parameters, $M_0$/$M_1$/$M_2$, $K$, $C$, $P$, $\fs$, $\ts$,
  \texttt{build\_M()}.
\item[lfr\_forward.py] 7-step resolve-and-retain forward pass.
\item[lfr\_simulate.py] RK4 step function, \texttt{simulate()} loop, \texttt{SimResult}.
\item[lfr\_block.py] \texttt{LFRBaselineBlock}---Jan's \texttt{Block} interface wrapper.
\item[validate\_lfr.py] All 5 validation sections (trajectory, Bode, Y-varying, eigenvalues,
  LPV vs frozen).
\item[test\_jan\_compat.py] Interconnect wiring tests and augmentation compatibility checks.
\end{description}

\end{document}
